\documentclass[12pt,a4paper]{article}
\usepackage[T1]{fontenc}
\usepackage[utf8]{inputenc}
\usepackage[margin=25mm]{geometry}
\usepackage{amsmath,amssymb,amsthm}
\usepackage{booktabs,array}
\usepackage{placeins}
\usepackage{xcolor}
\usepackage{graphicx}
\usepackage[super,sort&compress,comma]{natbib}
\usepackage{setspace}
\doublespacing
\usepackage[mathlines]{lineno}
\usepackage{etoolbox}
\makeatletter
\@ifpackagelater{lineno}{2022/07/29}{}{%
  \newcount\LN@ams@eqpen
  \cspreto{math@cr@}{%
    \global\@eqpen\numexpr\@eqpen+\LN@ams@eqpen\relax}
  \cspreto{align}{%
    \linenomath
    \postdisplaypenalty=0\relax
    \global\LN@ams@eqpen\interdisplaylinepenalty}
  \csappto{endalign}{%
    \global\LN@ams@eqpen\z@
    \endlinenomath}
}
\makeatother
\renewcommand{\linenumberfont}{\normalfont\sffamily\fontsize{8}{9.6}\selectfont\color{gray}}
\setlength{\linenumbersep}{12pt}
\modulolinenumbers[1]
\leftlinenumbers
\usepackage[font=small,labelfont=bf]{caption}
\usepackage[hidelinks]{hyperref}
\hypersetup{pdfauthor={Zhenglin Li, Peizhong Li, Ting Zhu and Yi Li}}
\urlstyle{same}
\renewcommand{\bibnumfmt}[1]{#1.}
\setlength{\bibsep}{0pt}
\raggedbottom
\newtheorem{pviproposition}{Proposition}
\setlength{\emergencystretch}{2em}
\setlength{\parskip}{0em}
\title{Supplemental material\\[0.4em]\large Physics-guided payload identification from aggregate bus power}
\author{Zhenglin Li, Peizhong Li, Ting Zhu and Yi Li\thanks{Corresponding author: Yi Li, Wuhan Institute of Technology, No. 206 Guanggu Street, Hongshan District, Wuhan 430205, China. Email: \texttt{13323107655@163.com}. ORCID: \href{https://orcid.org/0009-0000-9541-8300}{0009-0000-9541-8300}.}}
\date{}

\begin{document}
\maketitle
\linenumbers
\setcounter{section}{0}
\setcounter{equation}{0}
\setcounter{table}{0}
\renewcommand{\thesection}{S\arabic{section}}
\renewcommand{\theequation}{S\arabic{equation}}
\renewcommand{\thetable}{S\arabic{table}}

\section{Identifiability and domain bounds}
\label{app:pvi_certificate}

\begin{proof}[Proof of the conditional-injectivity proposition in the main article]
Because each component of $F$ is quadratic,
\begin{equation}
F(\beta_1)-F(\beta_2)=
J_\beta\!\left(\frac{\beta_1+\beta_2}{2}\right)(\beta_1-\beta_2).
\label{eq:pvi_midpoint}
\end{equation}
Convexity places the midpoint in $\Omega$. Full column rank implies that
the output difference vanishes only when $\beta_1=\beta_2$. Positive mass
then gives $c_x=h_x/m$ and $c_z=h_z/m$. Substituting
$\beta_1-\beta_2=S S^{-1}(\beta_1-\beta_2)$ and applying the singular-value
bound proves the stability inequality in the main article.
\end{proof}

To assess rank between grid points, we divide the physical solver box into
sub-boxes and enclose their mass-moment ranges. Since $J_\beta$ is affine,
$G_j=\partial J_\beta/\partial\beta_j$ is constant. For box center $b$ and
half-widths $r_j$, Mirsky's singular-value perturbation bound gives
\begin{align}
\ell&=\sigma_{\min}(J_\beta(b)S)
-(1+\varepsilon_G)\sum_j r_j\|G_jS\|_2-\delta,
\label{eq:pvi_box_bound}\\
\delta&=\varepsilon_{\mathrm{cover}}\max\!\left\{
1,\|J_\beta(b)S\|_2+(1+\varepsilon_G)\sum_jr_j\|G_jS\|_2\right\}.
\label{eq:cover_numerical_margin}
\end{align}
Here $\varepsilon_G=10^{-12}$ and $\varepsilon_{\mathrm{cover}}=10^{-10}$
are numerical safeguards. The diagonal entries of $S$ are the enclosing
domain spans: $0.19$~kg, $0.0476$~kg\,m, and $0.0125$~kg\,m. The unit floor
applies in these scaled coordinates. Bisection resolves all boxes. The
six-pose design has 17,797 leaves and a minimum bound of
$1.59647\times10^{-5}$; the eight-pose design has 3,327 leaves and a minimum
bound of $1.60902\times10^{-5}$. Because the partitions differ, these loose
bounds do not show which design is statistically more precise. Evaluation uses double precision
without directed outward rounding. The bounds therefore support the rank
condition numerically, without a formal rounding-error certificate.

Across an $11^3$ grid of the payload domain, repeated-pose,
constant-tool-pitch, and prescribed varying-pitch designs have ranks one,
two, and three, respectively. The eight-pose inverse recovered the noiseless
solution in all 675 solves over 27 objects and 25 initializations.
Independently calculated torques agree with MATLAB \texttt{rigidBodyTree}
within $6.67\times10^{-16}$~N\,m. These results are consistent with the
fixed-pose, nominal-calibration analysis. Rank at achieved postures under
plant mismatch requires separate evaluation.


\section{Electrothermal dynamics and power decomposition}
\label{app:cross_model}

The additional model evolves joint angle, velocity, motor current,
controller integral state, temperature, and an energy ledger. With quantities
referred to the output shaft, each drive obeys
\begin{align}
 L_i\dot i_i &= v_i-d_i(i_i)-R_i i_i-K_i\dot q_i,
 &\tau_i&=K_i i_i,\\
 R_i &= R_{i,0}\rho_i[1+\alpha_T(T_i-T_a)],
 &d_i(i_i)&=v_{d,i}\tanh(i_i/i_s)+r_{b,i}i_i,\\
 C_{T,i}\dot T_i&=R_i i_i^2-h_i(T_i-T_a),
 &P_{\mathrm{bus}}&=P_a+\sum_i\max(v_i i_i,0).
 \label{eq:electrothermal}
\end{align}
The same SI coefficient $K_i$ relates current to torque and velocity to
back EMF. An ideal nonregenerative 12~V bus supplies the drives; negative
bridge power is dissipated locally. Cold resistances satisfy
$R_{i,0}/K_i^2=b_i$, where $b_i$ are the nominal holding-loss coefficients.
Both models thus assume the same cold static calibration.
$P_a$ includes a prescribed offset and low-frequency ripple. Nominal geometry,
outer-loop gains, and stall-torque ratings are retained. Position error,
velocity damping, and integral torque set a saturated current reference.
A current loop follows, with bus-voltage saturation. This differs from the
original speed-dependent torque-envelope controller. Stribeck friction,
electrical dynamics, bridge loss, and temperature add further deviations.
Table 3 in the main article lists the parameters and severity-dependent
departures. Vectors follow joint order; link vectors follow the three pitch
links. Each nonzero departure is multiplied by severity $s$ and its frozen
uniform nuisance factor in $[0.8,1.2]$, except the common coefficient
$\alpha_T$ and auxiliary terms. The paired quantities below share factors.
We prescribe these values to study the mechanisms; they are not identified
from a physical device.



Baseline and mechanical-only cases have $\alpha_T=0$ and no added electrical
departures; electrical-only cases have no added mechanical departures.
Velocity, current, and integral torque start at zero at the first pose.
The model code specifies saturation and anti-windup rules.
Each link's mass and COM changes share one factor. Coulomb and Stribeck
terms share axis factors, as do the two bridge-loss terms. Auxiliary ripple
has a frozen uniform phase. The initial temperature rise and
$\alpha_T$ both scale with severity. Their product depends quadratically on
$s$, so severity is not a linear loss scale. The model excludes
mechanical backlash, structural compliance, contact, and source-voltage sag.
Individual simulated currents are excluded from estimator inputs.

MATLAB \texttt{ode15s} integrates the dynamics with relative tolerance
$\varepsilon_r=2\times10^{-6}$, absolute tolerance
$\varepsilon_a=10^{-8}$, and maximum step $h_{\max}=5$~ms. We assessed
numerical sensitivity on one complete 22.6~s trajectory at combined
$s=1$, using the fixed first payload from the development set.
We divided both tolerances by four and halved $h_{\max}$, keeping
observation times and random draws fixed. Angle, current, temperature,
and plateau power changed by at most $1.60\times10^{-7}$~rad,
$5.52\times10^{-6}$~A, $9.41\times10^{-7}$~K, and
$5.69\times10^{-6}$~W, respectively. This single trajectory measures
local numerical sensitivity, not convergence over the full domain.

For a saved holding sample, let $P_0$ be nominal power at its true angle
and payload, and $\tau_{0i}$ the corresponding nominal torque. The exact
signed accounting identity is
\begin{align}
 P_{\mathrm{bus}}-P_0={}&(P_a-a_0)+\sum_i d_i(i_i)i_i
 +\sum_i(R_i-R_{i,0})i_i^2 \notag\\
 &+\sum_iR_{i,0}\left[i_i^2-(\tau_{0i}/K_i)^2\right]
 +P_{\mathrm{shaft}}+P_{\mathrm{brake}}+P_{\mathrm{mag}},
 \label{eq:power_ledger}
\end{align}
where $P_{\mathrm{shaft}}=\sum_iK_i i_i\dot q_i$,
$P_{\mathrm{brake}}=\sum_i\max(-v_i i_i,0)$, and
$P_{\mathrm{mag}}=\sum_iL_i i_i\dot i_i$.
The resistance term separates into cold-scale and thermal contributions.
Magnetic power follows directly from the circuit equation. We evaluate the
identity before averaging, at 31 samples per pose, eight poses per episode,
and 156 episodes. These 38,688 samples quantify algebraic closure. The
payload object remains the independent unit for statistical inference.

To quantify mismatch that payload changes can absorb locally, evaluate
the nominal power Jacobian $J$ at true payload parameters and observed mean
angles. Let $p_{0,k}=P_{\mathrm{nominal}}(\bar q_k,\theta_{\mathrm{true}})$
at the observed mean angle. This differs from the samplewise true-angle
power in Eq.~\eqref{eq:power_ledger}. For the original six fitting poses,
write $A=\Sigma^{-1}JS_\theta$ and $b=\Sigma^{-1}(\bar P-p_0)$, where
$S_\theta$ scales parameter spans. The matrix
$\Sigma=\operatorname{diag}[\max(s_{P,k},\sigma_{P,\min})]$ uses the original
per-pose sample SD and floor $\sigma_{P,\min}=0.04$~W.
The least-squares step $u=A^\dagger b$ gives the diagnostic payload change
$\Delta\theta=S_\theta u$ and absorbed fraction
$\chi=\|Au\|_2^2/\|b\|_2^2$. We also evaluate the resulting power change on
the two withheld poses. This local projection may leave the physical box.
It is neither an estimate available to the robot nor a Fisher-information
calculation for the changed plant.

\section{Implementation details}
These settings supplement the estimator and simulation design in the main
article. Model selection uses only training and validation objects. The
true-state diagnostics in Section S2 are excluded from estimator inputs
and fitting.

\subsection{Probe windows and physical inversion}
The estimator uses aggregate power, native encoders, commands, and
timestamps, excluding measured torques and simulator states. The eight
targets in main-article Table 1 give joint angles in degrees in the
reduced-model coordinates. They provide holding-power observations across
postures. The 22.6~s probe starts with a 1~s wait, then eight 1.5~s
transitions, each followed by a 1.2~s hold. The first transition repeats the
starting pose; subsequent transitions use the quintic blend
$10u^3-15u^4+6u^5$, where $u$ is normalized transition time. The final
window of each hold supplies mean position $q_k$, mean power $\bar P_k$,
and power standard deviation $s_{P,k}$. Stationarity requires
\begin{equation}
 \frac{\|q(t_b)-q(t_a)\|_2}{T_w}<v_{\mathrm{stat}},\qquad
 \max_i\left(\max_{[t_a,t_b]}q_i-\min_{[t_a,t_b]}q_i\right)<q_{\mathrm{stat}},
 \quad t_b-t_a=T_w.
 \label{eq:stationarity_thresholds}
\end{equation}
The thresholds $T_w=0.3$~s, $v_{\mathrm{stat}}=0.035$~rad/s, and
$q_{\mathrm{stat}}=0.012$~rad screen residual motion. All eight windows
must be complete and stationary. Power SD measures within-window variation,
not the mean's standard error. The physical inverse (the teacher) fits poses 2--7 and
withholds poses 1 and 8 for evaluation. The final residual estimators use
all eight poses.

The solver box is $[0.06,0.25]$~kg by $[0.04,0.20]$~m by
$[-0.025,0.025]$~m; weights use $\sigma_k=\max(s_{P,k},\sigma_{P,\min})$ with $\sigma_{P,\min}=0.04$~W.
The solver, \texttt{lsqnonlin}, uses an analytic Jacobian and seven starts
in normalized box coordinates:
\begin{equation}
 \mathcal U_0=\left\{\begin{aligned}
 &(u_c,u_c,u_c),(u_l,u_l,u_l),(u_h,u_h,u_h),\\
 &(u_l,u_h,u_h),(u_h,u_l,u_l),(u_c,u_l,u_h),(u_c,u_h,u_l)
 \end{aligned}\right\}.
 \label{eq:inverse_starts}
\end{equation}
The values $(u_l,u_c,u_h)=(0.15,0.50,0.85)$ define separated interior
starts. We retain the solution with the lowest finite objective.
Acceptance requires a positive exit flag, numerical rank three, and
\begin{equation}
 \begin{aligned}
 \mathrm{RMSE}_{\mathrm{fit}}&\leq\varepsilon_{P,\mathrm{fit}},&
 \mathrm{RMSE}_{\mathrm{check}}&\leq\varepsilon_{P,\mathrm{check}},\\
 \sigma_{\min}(J_s)&\geq\underline\sigma_J,&
 \operatorname{cond}(J_s)&\leq\kappa_{J,\max},
 \end{aligned}
 \label{eq:inverse_acceptance}
\end{equation}
where $J_s$ is the fitting Jacobian scaled by the solver-box spans and
$\sigma_k^{-1}$. Main-article Table 2 gives the reconstruction and
sensitivity thresholds. \texttt{physics6} denotes this teacher. \texttt{physics8} uses the
same settings to fit all eight poses, omitting the held-out residual
criterion. Finite rejected estimates remain available for diagnosis.
Teacher rejection or \texttt{physics8} failure alone does not exclude an otherwise
finite learning record.


\subsection{Features, scaling and model selection}
For sensitivity features, evaluate $J_\theta$ at $\theta_0$ and define
\begin{equation}
 G=D_\sigma^{-1}J_\theta\operatorname{diag}(s_\theta),\qquad
 r=D_\sigma^{-1}[\bar P-P(\theta_0;\eta)],\qquad
 v=(G^{\mathsf T}G+\lambda_G I)^{-1}G^{\mathsf T}r.
 \label{eq:pnn_features}
\end{equation}
Here $s_\theta=(0.14~\mathrm{kg},0.14~\mathrm m,0.036~\mathrm m)$ and
$\sigma_w=0.04$~W. The diagonal matrix $D_\sigma$ contains $\sigma_k$;
$\lambda_G=10^{-3}$ provides fixed numerical damping. The 52-feature
representation appends 24 entries of $G$, eight $s_{P,k}/\sigma_w$
values, three entries of $v$, three $\log_{10}(1+\sigma_j(G))$ values,
and $\log_{10}(1+\operatorname{cond}(G))$, with a machine-precision
denominator floor. These features use only nominal calibration and measured
summaries, excluding labels, hidden nuisance factors, and transient
simulator information.

\texttt{ridge13} and \texttt{ridgeAug} regress the physical residuals $\theta-\theta_0$
on 13 and 52 features, respectively. They use squared error, penalized
coefficients, and an unpenalized intercept. \texttt{interactionRidge} adds the
91 products $z_i z_j$, $i\leq j$, of the first 13 standardized inputs
to the 52 features, then standardizes all 143 predictors. Each ridge
model searches the penalty grid in main-article Table 2 independently.

For $n$ eligible training records, define residuals
$u_i=\theta_i-\theta_{0i}-R(x_i)$ and amplitudes
$d_{R,j}=\max\{\sqrt{n^{-1}\sum_i u_{ij}^2},d_{R,\min,j}\}$. Minibatches estimate
the empirical training loss. \texttt{ridgeSkipNN13} and \texttt{ridgeSkipNNaug} use the
selected \texttt{ridge13} and \texttt{ridgeAug} models as fixed branch $R$; \texttt{scaledNN13} uses
$R=0$. Means, SDs, and RMS amplitudes come from training data alone.
SDs below $10^{-10}$ are replaced by one. Floors of $10^{-5}$~kg for mass
and $10^{-5}$~m per COM coordinate prevent unstable amplitude scaling.
Zero output-layer weights and biases make the initial skip model reproduce
ridge exactly. Training fixes $\alpha=1$. Validation rescales the completed
ensembles over the declared grid, including the ridge fallback $\alpha=0$.

Each of the three neural families uses six architecture/penalty
combinations and three seeds per combination, giving 54 fits. A compact
NN13 control adds three fits on the same data. It uses eight tanh units,
domain-span response scales, penalty 0.001, and \texttt{fitrnet} LBFGS
with 800 iterations and validation patience 30. Each member of the three
scaled families retains the epoch with the lowest smooth physical
validation loss, including epoch zero. Predictions are averaged equally
across seeds in physical units before scoring or common clipping.

\texttt{ridge13} minimizes validation mass MAE/$t_m$ plus mean planar
distance/$t_c$. Other candidates minimize the largest of eight
condition/metric ratios to the selected \texttt{ridge13} model. Ties are broken by
their mean, the pooled physical score, and fixed order. The primary
estimator is selected from the three scaled neural families using the first
two criteria and fixed family order. This rule differs from epoch
retention. Loss, preprocessing, and search budget differ between ridge and
neural models, so the comparison evaluates complete estimators rather than
nonlinearity alone.


\subsection{Nominal plant, sampling and mismatch distributions}
The yaw--pitch--pitch--pitch rigid-body arm is simulated in MATLAB R2026a
with local PID control, friction, actuator limits, and 1~ms integration.
Pitch-link lengths are
$(0.23870809,\allowbreak 0.21600955,\allowbreak 0.06811302)$~m and
masses are $(0.11642973,\allowbreak 0.15987536,\allowbreak 0.00475299)$~kg;
the prescribed unloaded mass is 0.826~kg. Gravity is $9.81$~m/s$^2$.
Nominal auxiliary draw is 2~W, and the holding coefficients are
$3.74335616(1,0.5,1,1)$ in W/(N\,m)$^2$. These simulation priors require
calibration for hardware use. The source and configuration specify all
inertia and controller parameters.

Host commands and encoder samples follow a 50~ms cadence, independent of
transport delay; power is sampled every 10~ms. Command and encoder-return
delays are independently uniform over 10--15~ms, rounded upward to the
integration grid. Undelayed power samples are paired with the latest
available encoder value. Gaussian encoder noise has SD $10^{-4}$~rad
before 4096-count/revolution quantization; power-noise SD is 0.03~W per
sample. Packet loss and clock drift are excluded.

Independent Latin-hypercube designs sample
\begin{equation}
 \mathcal D_{\mathrm{gen}}=[m_{g,\min},m_{g,\max}]
 \times[c_{x,g,\min},c_{x,g,\max}]
 \times[c_{z,g,\min},c_{z,g,\max}].
 \label{eq:generation_domain}
\end{equation}
The intervals are $[0.08,0.22]$~kg, $[0.05,0.19]$~m, and
$[-0.018,0.018]$~m. A dense sub-box within a fixed
$(0.256,0.060,0.080)$~m envelope determines physically admissible inertia
from mass and COM. Learning uses steady holding summaries with this inertia
construction. Training, validation, and test sets contain 12,000, 2,000,
and 2,000 disjoint objects. Training and validation use one episode per
object, balanced across four condition families. Each test object
undergoes all four conditions, giving 8,000 test episodes and 22,000
episodes overall.

The nominal condition has no added mismatch. For mechanical mismatch,
each prescribed departure is independently multiplied by a uniform
factor over 0.5--1.5: link mass/inertia scale departures
$(0.025,-0.02,0.03)$, longitudinal COM shifts $(1,-1,1.5)$~mm,
friction scale departures $(0.15,0.35,0.25,0.40)$, and PID gain
departures $(-0.02,-0.06,-0.04,-0.08)$. Electrical mismatch adds
\begin{equation}
 \Delta P=\alpha_P\sum_i|\tau_i|+\gamma_P\sum_i\tau_i^4.
 \label{eq:electrical_mismatch}
\end{equation}
The base coefficients, $\alpha_P=0.03285784$~W/(N\,m) and
$\gamma_P=0.02100986$~W/(N\,m)$^4$, are independently multiplied by
factors from the same uniform distribution. At their base values, the
terms contribute equally, totaling 2\% of nominal mean static power for
the known $(0.15,0.12,0)$ reference payload in SI units. The combined
condition includes both mechanisms. Nuisance, sensor, and communication
draws are independent between episodes. Conditions are paired by object,
not noise realization, and parameters stay fixed within each probe.
Training and test use the same prescribed mismatch distributions.

All learning methods require finite features, teacher estimates, and
targets, with at least 100 examples per family in each fitting split.
Failed records are retained without replacement. The sampling plan, code,
and seeds are fixed before generation. Model selection uses only training
and validation data; estimator parameters are fixed before test generation.
These data are disjoint from the reference-assisted experiment. Split
seeds are 1410101--1410103; sensor, communication, and nuisance seeds are
respectively 14200000, 14400000, and 14500000 plus trial ID.


\subsection{Matched affine controls and numerical fitting}
An affine control replaces the nonlinear branch in the main article's
prediction equation. It retains the 13 inputs, fixed ridge skip,
training-only standardization, response amplitudes, physical loss, and
three-seed averaging. The fixed pair uses the selected 16-unit nonlinear
branch and a 42-parameter affine branch, both with $\lambda_N=10^{-4}$ and
$\alpha=0.5$. Adam settings and validation stopping are identical. Equal
penalty coefficients, however, need not regularize the two function classes
equally. The affine search uses six penalties from $10^{-7}$ to $10^{-2}$
and the same five correction strengths. We evaluate fixed and
validation-selected variants, including the ridge fallback, on 500
additional objects under all four original-generator conditions, yielding
2,000 complete trajectories. Descriptive 95\% intervals use 5,000 paired
object resamples without multiplicity adjustment.

To assess sensitivity to optimization, we also solve the affine training
objective in double precision, fixing the original single-rounded design
and constants. Analytical gradients and Hessians are compared with finite
differences. Damped Newton iterations solve each of the six penalties
from zero and from the three retained Adam solutions. Training-gradient
and decrement tolerances are $\varepsilon_g=10^{-7}$ and
$\varepsilon_d=10^{-11}$, respectively; the latter scales with
$\max(1,|\mathcal L|)$. These criteria measure numerical stationarity
without using validation performance. Solutions must agree across
initializations and remain consistent under single-precision inference.
One comparator retains the fixed penalty and correction strength; the
other uses the original validation rule. This post hoc analysis reuses
the 500-object and 12-object panels, holding the nonlinear model fixed.


\subsection{Reference-assisted adaptation}
The secondary adaptation study uses a compact residual ensemble of three
members, each with eight hidden units, initialized on nominal data. The
principal ridge-skip estimator is not updated online. Each block stores
predictions for four unknown objects, then probes one known empty state,
four training references, and two distinct validation references, in
that order. In kilograms and meters, the training $(m,c_x,c_z)$ vectors are
$(0.09,0.06,-0.014)$, $(0.21,0.07,0.014)$,
$(0.10,0.18,0.012)$, and $(0.20,0.17,-0.012)$;
validation uses $(0.12,0.10,-0.005)$ and $(0.18,0.145,0.007)$.
Unknown ground truth serves only simulation generation and scoring.
Validation-reference labels never enter fitting.

The empty probe estimates absolute auxiliary draw $a\in[a_{\min},a_{\max}]$
and a multiplier $\lambda\in[\lambda_{\min},\lambda_{\max}]$ on the original
nominal pitch-axis holding coefficients. The fixed plausibility bounds are
$[a_{\min},a_{\max}]=[0,6]$~W and
$[\lambda_{\min},\lambda_{\max}]=[0.3,2.5]$. Six poses fit the calibration;
two are withheld. A proposal requires a scaled-design condition number at
most $10^5$, fitting and held-out RMSE at most $0.12$~W, and held-out RMSE
deterioration at most $0.004$~W. Updates change the network, calibration,
or both. Every calibration change must pass this screen, and reference
features are recomputed for each candidate. Network refitting combines the
nominal cache with up to 32 current and historical training-reference
measurements, weighted 1 and 30, respectively. Each member has a limit of
300 iterations and validation patience 30. One failed fit rejects the
candidate. Usable training references enter the replay pool even when the
update is not deployed.

Deployment uses raw predictions on the two current validation references.
Let $\ell_r=\|\operatorname{diag}(s_\theta)^{-1}
(\hat\theta_r-\theta_r)\|_2^2/3$, $L=(\ell_1+\ell_2)/2$, and
$s_\theta=(0.14~\mathrm{kg},0.14~\mathrm m,0.036~\mathrm m)^{\mathsf T}$.
The improvement and nominal-retention conditions are
\begin{equation}
\begin{aligned}
 L_{\mathrm{new}}&\leq L_{\mathrm{old}}-
       \max(\varepsilon_u,\rho_u L_{\mathrm{old}}),\\
 L_{\mathrm{nom}}(\eta_0,W_{\mathrm{new}})&\leq
       \gamma_{\mathrm{nom}}L_{\mathrm{nom}}(\eta_0,W_0).
\end{aligned}
\label{eq:adaptation_gate}
\end{equation}
The worst-reference constraint is
$\max_r\ell_{r,\mathrm{new}}\leq\allowbreak
\gamma_w\max_r\ell_{r,\mathrm{old}}+\varepsilon_u$.
Component-MAE increases are limited to $\Delta_k$.
Here $\varepsilon_u=10^{-6}$ excludes negligible improvements,
$\rho_u=0.05$ requires a relative reduction, and
$\gamma_{\mathrm{nom}}=1.25$ limits nominal forgetting. The fixed
allowances $\Delta=(0.15~\mathrm g,0.5~\mathrm{mm},0.5~\mathrm{mm})$
and $\gamma_w=1.05$ permit small local changes but limit component and
worst-reference deterioration. Settings are chosen independently of
unknown-object outcomes and fixed before evaluation. Reference features
must be usable, and physical acceptance coverage must not decrease.
$L_{\mathrm{nom}}$ averages $\ell$ over the compact model's fixed nominal
validation set, using original-calibration inputs and each learner's
initial model. The passing candidate with the lowest reference loss is
deployed. If none passes, calibration and weights stay unchanged. Scoring
retains every finite raw unknown prediction, including rejected estimates.
Unknown-object performance is evaluated separately from reference selection.

Four fixed scenarios---stationary, electrical, mechanical, and
combined---each contain eight object/noise streams of four blocks. This
gives 1408 probes, including 512 unknown events from 128 distinct payloads
paired across scenarios. Drift follows prescribed powers of the cubic
smoothstep $z(\rho)=3\rho^2-2\rho^3$, where $\rho$ is normalized episode
progress, with axis-specific exponents. Parameters stay constant within
each probe. Electrical drift reaches 1.30 times the nominal pitch holding
losses and 2.20~W auxiliary draw. Mechanical drift reaches friction
multipliers $(1.40,1.80,1.60,2.00)$ and PID-gain multipliers
$(0.92,0.85,0.88,0.82)$. Command delay is fixed at 20~ms, with no
encoder-return transport delay or added communication jitter. The frozen,
empty-only, gated-calibration, and full variants require 0, 1, 3, and 7
reference probes per opportunity; seven probes take 158.2~s, excluding
handling. Drift reaches its endpoint at the final reference probe, leaving
no subsequent unknown-object evaluation for the last update.

Adaptive ridge uses identical recorded observations, ordering, reference
weights, replay pool, calibration proposals, and gate thresholds. Both
learners fit the domain-normalized target $(\theta-\theta_0)\oslash s_\theta$.
Ridge refits 13 features with penalty $0.001$ and an unpenalized intercept,
dividing weighted squared error by total weight. It retains offline
standardization and uses a deterministic closed-form fit. Neural refits
recompute standardization and stop on validation performance. Each learner
retains its initial nominal-loss baseline; identical gates may therefore
accept different updates. Descriptive paired intervals use 4000 complete
object/noise-stream resamples without multiplicity adjustment.

\end{document}
